\documentclass[10pt,a4paper]{amsart}
\usepackage[margin=19mm]{geometry}
\usepackage{amsmath,amssymb,mathtools}
\usepackage[scr=rsfs,cal=euler]{mathalfa}
\usepackage{xfrac}
\usepackage[unicode,hidelinks,bookmarks=false]{hyperref}
\newcommand{\ie}{i.e.\ }
\newcommand{\cf}{cf.\ }
\newcommand{\resp}{respectively\ }
\newcommand{\Subsection}{\subsection}
\providecommand{\nobreakdash}{\nobreak\hbox{-}}
\setlength{\emergencystretch}{2em}
\multlinegap0pt
\usepackage{amssymb}
\newtheorem{thm}{Theorem}
\newtheorem{cor}{Corollary}
\title{Projective Kernels and Replicability in Modular Function Theory}
\author{Hicham Saber, Abdellah Sebbar}
\date{Source: arXiv:2609.12838v1 (CC BY 4.0)}
\begin{document}
\maketitle

\noindent\emph{Source and licence.} Projective Kernels and Replicability in Modular Function Theory. Version arXiv:2609.12838v1 (CC BY 4.0), distributed by arXiv under the Creative Commons Attribution 4.0 International licence, http://creativecommons.org/licenses/by/4.0/, as recorded in the arXiv licence metadata for this exact version and shown on the abstract page https://arxiv.org/abs/2609.12838v1. Full licence notice: \texttt{licenses/arxiv-2609.12838-CC-BY-4.0.txt}.\par\smallskip

\noindent\textit{Editorial scope.} Counted unit: the article's thm thm:Grunsky-quasimodular (source0.tex lines 403--421). The article's complete original proof of it is delivered with it (source0.tex lines 423--506, one \texttt{proof} environment of 84 lines). No mathematics is written, repaired or re-derived: the statement and the proof are the article's own lines, in the article's own notation, and no retained line is substituted or re-flowed.  Retained uncounted as the source-local context the proof needs: lines 184--214; lines 252--269; lines 324--328; lines 391--393.  Nothing was omitted from any retained span, so the package carries no omission record.  No retained line contains a citation, so the wrapper carries no bibliography and no bibliography entry.  No retained line contains a figure environment, an image-inclusion command or drawing code, so no image is delivered and the package declares no asset dependency.  Editorial formatting only, in addition: an AMS \texttt{amsart} preamble that restates the article's own declarations and macros used by the retained lines, namely the environments \texttt{thm}, \texttt{cor} on separate counters, together with the standard packages those lines need.  The retained environments print in this document as Theorem 1, Corollary 1 in order of appearance, whereas the article prints the same items with the numbering of its own full document.  The complete native source of the article as distributed in the arXiv e-print for this exact version is bundled unchanged as \texttt{sources/210155/source0.tex}.
\begin{thm}\label{thm:transformation-law}
Let \(f\) be a nonconstant meromorphic function on \(\mathbb H\). The following are equivalent.
\begin{enumerate}
\item The function \(f\) is \(\rho\)-equivariant for some projective representation of \(\Gamma\).
\item The function \(S_2[f]\) is a meromorphic modular form of weight \(4\) for \(\Gamma\).
\item For every \(n\geq2\), the higher Schwarzian \(S_n[f]\) is a meromorphic quasimodular form of weight \(2n\) and depth at most \(n-2\) for \(\Gamma\).
\end{enumerate}
When these conditions hold, for every \(n\geq2\) and every
\[
\gamma=
\begin{pmatrix}
a&b\\
c&d
\end{pmatrix}
\in\Gamma,
\]
one has
\begin{equation}\label{eq:universal-transformation}
(cz+d)^{-2n}S_n[f](\gamma z)
=
\sum_{j=0}^{n-2}\binom{n-2}{j}S_{n-j}[f](z)
\left(\frac{c}{cz+d}\right)^j.
\end{equation}
The corresponding transformation polynomial is
\[
P_n(X)
=
\sum_{j=0}^{n-2}\binom{n-2}{j}S_{n-j}[f]X^j.
\]
Its degree is at most \(n-2\), and its coefficient of \(X^{n-2}\) is \(S_2[f]\). Whenever the associated quasimodular polynomial is unique, \(S_n[f]\) therefore has exact depth \(n-2\) if \(S_2[f]\not\equiv0\). If \(S_2[f]\equiv0\), then \(f\) is a linear fractional function and \(S_n[f]\equiv0\) for every \(n\geq2\).
\end{thm}

\begin{cor}\label{cor:lowering}
Suppose that $f$ is equivariant and the associated quasimodular polynomials for $\Gamma$ are unique. Then
\begin{equation*}
\delta S_n[f]=(n-2)S_{n-1}[f],
\quad n\geq3.
\end{equation*}
More generally,
\begin{equation}\label{eq:delta-iterate}
\delta^rS_n[f]
=
\frac{(n-2)!}{(n-2-r)!}S_{n-r}[f],
\quad 0\leq r\leq n-2.
\end{equation}
In particular,
\[
\delta^{n-2}S_n[f]=(n-2)!S_2[f].
\]
\end{cor}

\begin{equation}\label{eq:grunsky-kernel-expansion}
\mathcal G_f(q,r)
=
\sum_{m,n\geq1}mn\,h_{m,n}q^mr^n.
\end{equation}

\begin{equation}\label{eq:Bernoulli}
\frac{x}{e^x-1}=\sum_{j\geq0}B_j\frac{x^j}{j!}.
\end{equation}

\begin{thm}\label{thm:Grunsky-quasimodular}
For every $k\geq2$,
\begin{equation}\label{eq:moment-Schwarzian}
S_k[F](\tau)
=
-\frac{\kappa^k}{(k-1)!}\,\mathfrak G_k[F](\tau).
\end{equation}
Assume that $F$ extends meromorphically to $\mathbb H$. Let $\Gamma$ be a discrete subgroup of $\mathrm{SL}_2(\mathbb R)$ containing the translation $\tau\mapsto\tau+\ell$. Then:
\begin{enumerate}
\item $F$ is projectively equivariant for $\Gamma$ if and only if the local series $\mathfrak G_2[F]$ extends meromorphically to $\mathbb H$ as a modular form of weight $4$ for $\Gamma$;
\item if $F$ is projectively equivariant, then for every $k\geq2$ the local series $\mathfrak G_k[F]$ extends meromorphically to $\mathbb H$ as a quasimodular form of weight $2k$ and depth at most $k-2$; if the associated polynomial is unique and $S_2[F]\not\equiv0$, the depth is exactly $k-2$;
\item whenever the associated polynomials are unique, for $k\geq3$,
\begin{equation}\label{eq:Grunsky-lowering}
\delta\mathfrak G_k[F]
=
\frac{(k-1)(k-2)}{\kappa}\,\mathfrak G_{k-1}[F].
\end{equation}
\end{enumerate}
\end{thm}

\begin{proof}
For a locally univalent function $g$, set
\[
\mathcal P_g(u,v)
=
\frac{g'(u)g'(v)}{(g(u)-g(v))^2}-\frac1{(u-v)^2}.
\]
Compare the kernels for $F=f\circ q$ and for the local parameter $q$. Since $q'(\tau)=\kappa q(\tau)$,
\[
\frac{F'(\tau)F'(\sigma)}{(F(\tau)-F(\sigma))^2}
=\kappa^2q(\tau)q(\sigma)
 \frac{f'(q(\tau))f'(q(\sigma))}
 {(f(q(\tau))-f(q(\sigma)))^2}.
\]
Similarly,
\[
\mathcal P_q(\tau,\sigma)
=\kappa^2\frac{q(\tau)q(\sigma)}{(q(\tau)-q(\sigma))^2}
-\frac1{(\tau-\sigma)^2}.
\]
Subtracting these two expressions gives the local parameter cocycle
\begin{equation}\label{eq:kernel-cocycle}
\mathcal P_F(\tau,\sigma)
=
\mathcal P_q(\tau,\sigma)
-\kappa^2\mathcal G_f(q(\tau),q(\sigma)).
\end{equation}

Set $\sigma=\tau+t$. The defining expansion of the Aharonov invariants gives
\[
\mathcal P_F(\tau,\tau+t)
=\sum_{k\geq2}(k-1)S_k[F](\tau)t^{k-2}.
\]
For the exponential coordinate $q=e^{\kappa\tau}$,
\[
\frac{q'(\tau)}{q(\tau+t)-q(\tau)}
=\frac{\kappa}{e^{\kappa t}-1}.
\]
Comparison with the Bernoulli generating function \eqref{eq:Bernoulli} yields
\begin{equation*}
S_k[q](\tau)=-\frac{B_k\kappa^k}{k!},
\quad k\geq2.
\end{equation*}
The coefficient of $t^{k-2}$ in $\mathcal P_q(\tau,\tau+t)$ is
\[
(k-1)S_k[q]
=-\frac{B_k\kappa^k}{k(k-2)!}.
\]

The Grunsky term is expanded from \eqref{eq:grunsky-kernel-expansion}:
\[
\mathcal G_f(q,qe^{\kappa t})
=
\sum_{m,n\geq1}mn\,h_{m,n}q^{m+n}e^{\kappa nt}.
\]
Its $t^{k-2}$ coefficient is
\[
\frac{\kappa^{k-2}}{(k-2)!}
\sum_{m,n\geq1}m n^{k-1}h_{m,n}q^{m+n}.
\]
Taking the coefficient of $t^{k-2}$ in \eqref{eq:kernel-cocycle} and dividing by $k-1$ gives
\[
S_k[F]
=-\frac{\kappa^k}{(k-1)!}
\left(
\frac{B_k}{k}
+\sum_{m,n\geq1}m n^{k-1}h_{m,n}q^{m+n}
\right),
\]
which is \eqref{eq:moment-Schwarzian}.

For $k=2$, this identity expresses $S_2[F]=\{F,\tau\}/6$ as a nonzero constant multiple of $\mathfrak G_2[F]$. The Schwarzian criterion gives assertion (1): modular continuation of $\mathfrak G_2$ is equivalent to projective equivariance of $F$. If $F$ is equivariant, Theorem~\ref{thm:transformation-law} gives the transformation polynomial of each $S_k[F]$. Multiplication by the constant in \eqref{eq:moment-Schwarzian} gives the corresponding meromorphic continuation and quasimodular transformation law for $\mathfrak G_k[F]$, which is assertion (2).

Assume now that the associated quasimodular polynomials are unique. Applying $\delta$ to \eqref{eq:moment-Schwarzian} and using
\[
\delta S_k[F]=(k-2)S_{k-1}[F]
\]
from Corollary~\ref{cor:lowering} gives
\[
-\frac{\kappa^k}{(k-1)!}\delta\mathfrak G_k
=-(k-2)\frac{\kappa^{k-1}}{(k-2)!}\mathfrak G_{k-1}.
\]
Cancelling the common factors gives \eqref{eq:Grunsky-lowering}.
\end{proof}

\end{document}
